% Incompressible Euler Flow: Turbulence vs Blow-up
% Johan Hoffman and Claes Johnson. DRAFT.
\documentclass[11pt]{article}
\usepackage[margin=1.15in]{geometry}
\usepackage{amsmath,amssymb}
\usepackage{url}
\usepackage[hidelinks]{hyperref}
\setlength{\parskip}{3pt}

\newcommand{\grad}{\nabla}
\newcommand{\eps}{\varepsilon}

\title{\vspace{-2em}Incompressible Euler Flow: Turbulence vs Blow-up}
\author{Johan Hoffman\\[0.2em]
\small KTH Royal Institute of Technology, Stockholm
\and Claes Johnson\\[0.2em]
\small KTH Royal Institute of Technology, Stockholm}
\date{\today}

\begin{document}
\maketitle

\begin{abstract}
\noindent
All physical Euler flow is a form of bluff body flow. Without an obstruction in the flow, and with
slip on the boundary, nothing happens: a uniform stream in a straight pipe solves the equations for
ever, there being no shear to be unstable and no viscosity to make a wall layer. Euler with slip has
no mechanism for generating structure, so everything that makes a flow a flow enters because the
fluid must go around something. We therefore take flow past a body as the only setting in which the
equations describe a fluid. There potential flow is a solution and it is unstable, but not
everywhere. Attachment, on the forward surface where the flow
is accelerated, is stable. Separation, on the rear surface where it is retarded, is not: the
opposing flow there is unstable to spanwise perturbation and generates streamwise vorticity, which
stretches, breaks down and feeds a cascade. The flow that results has bounded velocity and
unbounded velocity gradients, is H\"older continuous with exponent about $1/3$, and dissipates
energy at a rate that does not vanish --- a dissipative weak solution in the sense of Onsager. This
is what flow past a body does, and the distinction between streamlined and bluff bodies does not
survive it: every body separates, so all flow past a body is bluff body flow, and there is no
potential flow regime at any Reynolds number.

We then compare this with the programme that constructs finite-time singularities of Euler, which
has succeeded in several settings. Our conclusion is that \emph{all physical Euler flow is
turbulent and the blow-up solutions are not physical}. They are theorems and we do not dispute
them; they are not flows. The reason is a matter of function space. The singular solutions are
constructed in $C^{1,\alpha}$, far above Onsager's exponent $1/3$, where energy is necessarily
conserved; a solution that conserves energy has no cascade available to it, and nowhere to put a
growing gradient but concentration. Physical flow lies below the exponent, where dissipation is
available, and takes it. Blow-up is what Euler does when it cannot dissipate; turbulence is what it
does when it can, and a fluid always can. The constructions are moreover posed either on the whole
space, where there is no body and so no separation, or with a symmetry and a boundary that supply
the concentration by hand --- and stability within a symmetry class is not genericity, least of all
when the physical instability is the one that breaks the symmetry.
\end{abstract}

\section{All physical Euler flow is bluff body flow}\label{sec:setting}

We begin with the claim the title rests on, because it is not a summary of what follows but the
reason there is anything to follow.

\emph{All physical Euler flow is a form of bluff body flow.} Without an obstruction in the flow, and
with slip on the boundary, nothing happens at all. Take a straight pipe and a uniform stream along
it. With $\nu=0$ and slip on the wall, $u$ constant and $p$ constant solve the Euler equations exactly,
and go on solving them for ever. There is no shear, so there is nothing to be unstable; there is no
wall layer, because there is no viscosity to make one; and there is nothing for the fluid to flow
past. The solution is the initial data, unchanged, for all time. One may compute it, but there is
nothing in it to compute.

This is not a curiosity about pipes, and the general statement should be made exactly, because it is
the hinge of the paper. \emph{The Euler equations do not create vorticity.} The vorticity equation
$\dot\omega + (u\cdot\grad)\omega = (\omega\cdot\grad)u$ transports and stretches $\omega$ and has no
source term; with slip and no viscosity there is none at the wall either. So the vorticity in an
Euler flow is either present in the initial data, or it is produced where the fluid is obliged to go
around something.

That is a dichotomy, not a preference, and it divides the subject. If the vorticity is in the data,
it was put there by whoever wrote the data down, and what is being studied is an initial-value
problem: the flow's structure is an input. If it is produced at an obstruction, the flow generates
its own structure and one is studying a fluid. Everything that makes a flow a flow --- vorticity,
separation, a wake, drag, turbulence --- belongs to the second case, and no fluid ever arrives at the
first, because a fluid acquires its vorticity from boundaries. The obstruction is not one case among
many. It is the only physical source.

And once there is an obstruction, there is a forward surface where the flow is accelerated and a
rear surface where it is retarded, because the pressure must return to what it was. The first is
stable and the second is not, as Section~\ref{sec:dalembert} sets out, and from the instability of
the second everything else follows. That is bluff body flow. A streamlined body retards the flow more
gently over a longer rear surface and so separates later and leaves a narrower wake, but it
separates, and the mechanism is the same one. There is no third kind of body and no flow without a
body worth the name. Hence the claim, and hence also the shape of this paper: the setting is not a
special case chosen for convenience, it is the only setting in which the Euler equations describe
anything.

\subsection{Euler with slip, past a body}\label{sec:eqns}

The incompressible Euler equations
\begin{equation}\label{eq:euler}
\dot u + (u\cdot\grad)u + \grad p = 0,\qquad \grad\cdot u = 0,
\end{equation}
are posed in a domain $\Omega$ exterior to a body $B$, with $u\to U$ at infinity and the
\emph{slip} condition
\begin{equation}\label{eq:slip}
u\cdot n = 0 \quad\text{on}\quad \partial B,
\end{equation}
no condition being imposed on the tangential velocity.

Neither of these is a free choice, and both are worth a word, because the first is routinely
treated as one and the second is usually not noticed at all.

\paragraph{Why slip.} Because the viscosity is zero. The slip
condition~\eqref{eq:slip} is not a modelling choice to be weighed against no-slip; it is the only
boundary condition the Euler equations admit on an impermeable wall. The wall must not be crossed,
which fixes $u\cdot n$; nothing in~\eqref{eq:euler} determines the tangential velocity there, and
imposing $u=0$ would over-determine the problem. No-slip belongs to Navier--Stokes, where the
viscous term is second order and carries the extra condition, and the boundary layer it produces has
thickness $\sim Re^{-1/2}$ and disappears in the inviscid limit. At $\nu=0$ there is no boundary
layer to be thin. There is slip.

This should be said because the drag of a body is habitually attributed to the boundary layer, and
in the present setting there is none. Whatever drag appears below is produced without viscosity and
without a wall layer, by the instability of Section~\ref{sec:dalembert} alone. That the resulting
computations give the right numbers --- mean drag within a few per cent for the circular cylinder,
the separation line in the observed place, the lift curve and the stall for
wings~\cite{drag,flight} --- is then evidence that nothing essential was left out. The
turbulent Euler solutions obtained this way on coarse meshes, including the drag crisis, are
reported in~\cite{dragcrisis}.

\paragraph{Why past a body.} The Cauchy problem for~\eqref{eq:euler} on $\mathbb{R}^3$ or on the
torus is a perfectly well-posed mathematical question and we do not suggest otherwise. We observe
only that it is not a question about a fluid. A fluid flows past something; with no boundary there
is no attachment and no separation, hence none of the mechanism described below, and what is being
studied is an initial-value problem for a quasilinear system rather than a flow. This matters later,
because several of the singularity constructions are posed exactly there.

\section{d'Alembert: attachment is stable, separation is not}\label{sec:dalembert}

Potential flow --- the irrotational solution of~\eqref{eq:euler} with~\eqref{eq:slip} --- exists for
any body, is unique, and exerts zero drag. That is d'Alembert's paradox, and for a century and a
half the standard resolution located the missing drag in viscosity and the boundary layer. The
resolution we gave in~\cite{dalembert} locates it elsewhere: potential flow is a solution of the
equations, and it is \emph{unstable}, so it is not observed. Nothing about viscosity is required to
say so.

What matters for everything that follows is that the instability is not uniform over the body.
Divide the surface by the sign of the streamwise pressure gradient.

On the forward surface the flow is \emph{accelerated}: the pressure gradient is favourable, fluid
elements are stretched along the flow direction and compressed across it, and a perturbation to the
velocity is correspondingly stretched and weakened. Attachment is stable. This is why the forward
stagnation point and the flow over the front of a body are seen in experiment exactly where
potential theory puts them, and why potential flow is an excellent description of the forward
surface.

On the rear surface the flow is \emph{retarded}: the pressure gradient opposes the motion, and
potential flow asks opposing streams to meet at a rear stagnation point and separate smoothly. That
configuration is unstable. Perturbing it in the spanwise direction, the opposing flows do not meet
on a line but interleave, and the result is a row of counter-rotating vortex tubes aligned with the
flow. The growth is exponential with rate set by $|\grad u|$ at separation, so it is fast on the
convective time scale, and it is not a slow accumulation that a sufficiently careful experiment
could avoid.

The asymmetry is the whole of the resolution. The drag is zero in potential flow because the
pressure recovers on the rear exactly as it fell on the front. The instability destroys the recovery
and only the recovery: the front is as potential theory says, the rear is not, and the difference is
the drag.

\section{What the instability generates}\label{sec:vorticity}

The perturbation that grows at separation is specific, and its form is what connects a local
instability to a global cascade.

The unstable mode generates \emph{streamwise} vorticity: vortex tubes aligned with the flow, in
counter-rotating pairs, originating at the separation line. This is not incidental. A vortex tube
aligned with the flow sits in the strain field of the flow that produced it, and that field stretches
it. Stretching a vortex tube intensifies its vorticity and thins its core, which is to say it moves
energy from the scale of the body to a smaller scale, and it does so without any help. The stretched
tube is then itself unstable, breaks down into finer tubes, and the process repeats.

What this gives is a cascade with an origin. In the usual presentation the Kolmogorov cascade is a
statistical hypothesis about an inertial range, with the large scales stirred by some unspecified
mechanism. Here the mechanism is named and located: it is the separation instability, it acts at the
rear of the body, and the first generation of the cascade is the counter-rotating streamwise tubes
it produces. The inertial range is what happens to them.

Energy enters this process at a rate set by the body and the free stream, $\eps\sim U^3/L$, and
nothing in the chain depends on viscosity. Each generation passes its energy down, and the rate is
the same at every scale because there is nowhere else for it to go.

\section{No streamlined exception}\label{sec:bluff}

Section~\ref{sec:setting} claimed that all physical Euler flow is bluff body flow, and the
instability of Section~\ref{sec:dalembert} is what makes the claim bite. The objection to meet is
the conventional division of bodies into bluff and streamlined, the first separating and leaving a
wide wake, the second supposedly not. The division does not survive the account above.

Every body separates, because every body is finite and the flow that accelerates over its forward
surface must be retarded over its rear. A streamlined body is one whose rear surface retards the
flow gently and whose separation is therefore delayed and its wake narrow, and this is a difference
of degree in the size of the unstable region, not a difference in kind. There is no shape for which
the rear pressure gradient is favourable everywhere, because the pressure must return to its
free-stream value.

\paragraph{This is a record, not an assertion.} The statements above are not predictions about
what would happen to a perturbed potential solution. Taking $\nu\to0$ in the General Galerkin method
of~\cite{book} gives a solver for~\eqref{eq:euler} whose dissipation does not vanish with the
viscosity, and what it produces, started from smooth data and with no model imposed to keep it
smooth, is not potential flow. A bluff body in this limit shows potential flow modified by a
three-dimensional rotational separation, and that separation carries the drag. The turbulent Euler
solutions so obtained, including the drag crisis, are reported in~\cite{dragcrisis}; the same
computation applied to a wing gives lift and drag with no circulation condition imposed by
hand~\cite{flight}.

Seen this way d'Alembert's paradox is a mistaken limit rather than a missing force. It arises from
taking the smooth potential solution to be the vanishing-viscosity limit of a real flow. The correct
limit is turbulent, and it has drag.

Two consequences follow and both are testable. There is no potential flow regime: potential flow is
unstable past every body at every Reynolds number, so the limit $Re\to\infty$ does not recover it,
and the inviscid limit of a real flow is turbulent and not smooth. And the same mechanism accounts
for a wing: a wing is a slender bluff body whose separation is organised rather than suppressed, and
the lift it generates comes about because the instability acts at the trailing edge and on the upper
surface in a way the geometry controls. That is the account of flight given in~\cite{flight}, and it
is the same instability as the cylinder's.

\section{What the physical solution is}\label{sec:weak}

The flow produced by this process is not a classical solution of~\eqref{eq:euler}, and it is worth
being exact about what it is instead, since the comparison in Section~\ref{sec:blowup} turns on it.

Energy is bounded: the velocity is bounded by the free stream and the domain is fixed, so
$\|u(t)\|_{L^2}$ does not grow. Velocity \emph{gradients} are not bounded. The cascade refines
without limit, since there is no viscosity to stop it, and the velocity difference across a
separation $\ell$ behaves as
\begin{equation}\label{eq:holder}
\delta u(\ell)\sim(\eps\ell)^{1/3},
\qquad\text{so}\qquad
\frac{\delta u(\ell)}{\ell}\sim \eps^{1/3}\ell^{-2/3}\to\infty .
\end{equation}
The first relation is Kolmogorov's~\cite{kolmogorov41}, and read as a statement about the velocity
field it says that $u$ is H\"older continuous with exponent $1/3$ and no better.

That exponent is not a detail. Onsager observed~\cite{onsager} that a weak solution of Euler
conserves energy if its velocity is H\"older continuous with exponent greater than $1/3$, and that
below the exponent it need not. The first half was proved by Constantin, E and Titi~\cite{cet}; the
second half --- that dissipative weak solutions below the exponent exist --- was established through
the convex-integration programme of De Lellis and Sz\'ekelyhidi and
collaborators~\cite{dls,bdlis} and completed by Isett~\cite{isett}. Eyink's analysis~\cite{eyink}
identifies the mechanism as a local energy flux that survives the inviscid limit.

Onsager's threshold is now settled in both directions --- above the exponent energy is conserved,
below it dissipating solutions exist and have been constructed --- so the class we are describing,
bounded and not smooth and dissipating, is not a conjecture about what solutions might do. It is a
category with occupants.

So the physical solution is a \emph{dissipative weak solution}: bounded velocity, unbounded
gradient, H\"older exponent at the Onsager value, and energy dissipated at a rate $\eps$ that does
not vanish although the equation contains no dissipative term. The dissipation is anomalous in the
technical sense and entirely ordinary in the physical one --- it is the drag, and it has been
measured since Taylor~\cite{taylor}.

That the dissipation in external flow past a body is inertial rather than viscous, and survives
the inviscid limit, is not only our claim: it has been pursued independently within the Onsager
framework for flow around a smooth body~\cite{inertialbody}, where the momentum dissipation is
located at the boundary and in the wake. We take the agreement as support for the mechanism and note
the difference in route --- that work argues from the viscous side and takes a limit, and the
account here starts at $\nu=0$ and never has a boundary layer to dispose of.

\paragraph{What is proved, and what is computed.} The three statements above do not have the same
standing and it is better to separate them than to let a reader assume they do.

\emph{That the limit is a weak solution of Euler is proved.} General Galerkin is a stabilised
finite element method whose stabilisation is a weighted least-squares control of the residual: the
energy estimate for the discrete problem bounds the residual of the computed solution in a weighted
$L_2$ norm, and that bound is what the method buys in place of the smoothness it does not assume.
Testing the residual against a smooth function and using the weighting, the residual tends to zero
in the sense required by the weak formulation, so the limit of the computed sequence satisfies
Euler's equations weakly~\cite{book}. This is not an identification offered on the strength of the
numbers looking right. It is a theorem, and it is the reason G2 is a method for Euler rather than a
method for some nearby regularised equation.

\emph{That the dissipation does not vanish is computed.} The rate is measured, and measured to be
independent of the mesh and of the vanishing viscosity, which is what makes the weak solution a
\emph{dissipative} one. We know of no proof of this and do not claim one.

\emph{That the exponent is exactly $1/3$ is inferred.} Relation~\eqref{eq:holder} is Kolmogorov's,
read as a statement about the velocity field, and we do not derive it
from~\eqref{eq:euler}. What we claim is the mechanism that puts a flow into that class and keeps it
there: separation is unstable, the instability generates streamwise vorticity, stretching carries it
down the scales, and the cascade terminates nowhere because nothing terminates it. Onsager's
theorem then says that a dissipative weak solution \emph{cannot} have exponent above $1/3$, so the
proved and computed halves together bound the exponent from one side without the phenomenology.

\section{The constructed singular solutions}\label{sec:blowup}

Against this stands a programme, now substantially successful, that constructs solutions
of~\eqref{eq:euler} whose \emph{gradient} becomes infinite in finite time. Luo and Hou exhibited a
numerical scenario with boundary~\cite{luohou}; Elgindi proved finite-time singularity formation for
axisymmetric flow without swirl at $C^{1,\alpha}$ regularity~\cite{elgindi}; Chen and Hou gave a
computer-assisted proof for the 2D Boussinesq and 3D axisymmetric Euler equations with boundary and
smooth initial data~\cite{chenhou}; and the regularity has since been pushed to
$C^{1,1/3}$~\cite{chen2026}, with stability on the whole space also claimed~\cite{stable}.

\paragraph{What is actually proved, and it is not what the name suggests.} The word ``blowup''
invites the reading that the velocity becomes infinite. It does not, and the theorems do not claim
it. Taking the most recent and strongest, Theorem~1.1 of~\cite{openaieuler} asserts that there is
smooth compactly supported divergence-free $u_0$ whose smooth Euler solution has finite maximal
lifespan $T_*$, with
\[
\limsup_{t\uparrow T_*}\|\grad u(t)\|_{L^\infty}=\infty,
\qquad
\int_0^{T_*}\|\operatorname{curl}u(t)\|_{L^\infty}\,dt=\infty .
\]
What diverges is the \emph{gradient}, and the vorticity in the Beale--Kato--Majda
integral~\cite{bkm}. Nothing is asserted about the velocity becoming unbounded.

That is the same signature as the turbulent solution of Section~\ref{sec:weak}: gradients diverge,
velocity need not. So the disagreement is narrower than the names suggest, and it is worth saying so
plainly rather than being shown it. Both accounts hold that a smooth Euler solution ends in finite
time with unbounded gradient. We said as much computationally in~\cite{blowupHJ}, and it is the
content of Section~\ref{sec:dalembert} here. Nor is the mechanism as foreign as one might expect:
the construction of~\cite{openaieuler} proceeds by adding a localised oscillation to a parent flow
at each stage, amplifying it, and iterating over scales to a finite limiting time --- a cascade,
built rather than generated.

The disagreement is therefore not about whether smooth solutions break down. It is about two other
things: where the vorticity came from, and whether energy leaves. On both the answers are opposite,
and they are what the rest of this section is about.

\paragraph{The unforced construction of September 2026.} The sharpest case is the most recent. On
26 September 2026 OpenAI released a proof of finite-time blowup for the three-dimensional
incompressible Euler equations, \emph{unforced}, from smooth, compactly supported, divergence-free
initial data, accompanied by a formalisation in Lean~\cite{openaieuler}. Comparable results on Euler were being
prepared independently and by hand at the same time~\cite{buckmaster}, so what follows is not an
argument about machines: the construction is a real object in the subject and would be there without
them. It is the strongest form of
the claim: no forcing to object to, no low-regularity data, and a machine-checked proof. We take it
as correct and our argument does not depend on whether it is --- what follows would be unaffected by
a verification or by a retraction.

It is also the cleanest illustration of both points above, which is why we give it separately.

Its data is \emph{smooth}, and that settles the matter without appeal to any delicate theorem. Local
well-posedness for Euler with smooth, compactly supported data gives a unique \emph{classical}
solution with persistence of regularity on a maximal interval $[0,T)$, and by Beale--Kato--Majda the
interval is finite exactly when the vorticity ceases to be integrable in time~\cite{bkm}. So
\emph{blowup at $T$} means precisely that the solution is smooth on all of $[0,T)$. For a classical
solution with decaying data,
\[
\frac{d}{dt}\int \tfrac12|u|^2
= -\int \grad\cdot\Big(u\tfrac{|u|^2}{2}\Big) - \int \grad\cdot(pu) = 0 ,
\]
an integration by parts and nothing more: the energy is conserved exactly on the whole interval of
existence. The construction therefore has no cascade at any moment of its existence, and the sharpest way to
put that is the plainest: \emph{it dissipates nothing}. Not a little, not slowly --- nothing, over
the whole of $[0,T)$. The singularity is accordingly not a mechanism by which energy is lost. It is
a mechanism by which energy is \emph{piled up}: at the singular time every bit of it is still
present, gathered into a vanishing region.

Set against that the flow of Sections~\ref{sec:dalembert}--\ref{sec:weak}, which begins losing
energy as soon as it meets the body and goes on losing it at a rate that does not depend on any
viscosity --- a loss which is the drag, and has been measured since Taylor~\cite{taylor}.

\paragraph{The construction is d'Alembert's paradox in time-dependent form.} This is worth stating
as the unification it is. Potential flow is smooth, conserves energy, and has zero drag: that is the
paradox of Section~\ref{sec:dalembert}. The blowup solution is smooth up to $T$, conserves energy,
and dissipates nothing. The turbulent solution is not smooth, dissipates, and has drag. The first
two are what one obtains by insisting on smoothness, and they fail as descriptions of a fluid for
one and the same reason: a fluid dissipates, and a smooth Euler solution provably cannot. The
resolution is also the same resolution. The smooth solution exists, and is not the one that occurs. Its energy cannot leave, and as the gradient
grows the integral can only be kept by concentrating. The singularity is not something the flow
discovers; it is the only remaining option, and the smoothness is what removed the others.

And its data is compactly supported on $\mathbb{R}^3$ --- so there is no body, no wall and no
separation, and by Section~\ref{sec:setting} the vorticity it starts with was put there by hand. It
is a divergence-free field chosen so that the stretching term concentrates it, and the proof is the
demonstration that it does. No fluid arrives at that state. A fluid gets its vorticity from a
boundary, which is to say from the separation instability, which is to say with the roughness that
instability immediately produces --- and a rough flow has the cascade available and takes it. The
construction and the flow differ at the first step, in where the vorticity came from, and everything
else follows from that.

We have been here before, and it is worth saying where we stood. In~\cite{blowupHJ} we studied
blowup of incompressible Euler solutions computationally and concluded that what the equations
produce is not a point singularity but a breakdown of the flow into turbulence --- the velocity
staying bounded while the gradients do not. The present paper is that conclusion stated in the
language the subject has since acquired, with Onsager's exponent as the thing that separates the two
outcomes, and with the separation instability named as the mechanism that drives a flow across it.

These are theorems and we do not dispute them. We claim that they are not flows. All physical
Euler flow is turbulent, by the mechanism of Sections~\ref{sec:dalembert}--\ref{sec:vorticity},
which acts on every body and has no exceptions; the singular solutions are solutions of the
equations that no fluid realises. The reason is not rhetorical but a matter of which function space
each account occupies, and it can be stated exactly.

\paragraph{The regularity is the point.} Every one of these constructions lives in $C^{1,\alpha}$:
the velocity is continuously differentiable with H\"older continuous derivatives. That is far
\emph{above} Onsager's exponent. A $C^{1,\alpha}$ velocity is Lipschitz, so these are classical
solutions up to the singular time, and conserve energy exactly there by the same integration by
parts. Constantin, E and Titi's theorem~\cite{cet} is not needed for this; it is the sharp statement
for \emph{weak} solutions, and it is what fixes $1/3$ as the threshold below which conservation can
fail. What matters here is only that every one of these constructions sits far above it.

A solution that conserves energy has no cascade available to it. The energy present cannot be
passed down the scales and removed; it must remain, and the only way to accommodate a growing
gradient while conserving the integral is to concentrate. The singularity is what that
concentration becomes --- it is forced, and it is forced by the smoothness. That is the sense in
which these solutions are not physical. A fluid is not confined to $C^{1,\alpha}$. It is free to
roughen, the separation instability makes it do so immediately, and once it is below the Onsager
exponent the cascade is open and the energy leaves. The concentration never has to happen, so it
does not.

\begin{quote}
Blow-up is what Euler does when it cannot dissipate. Turbulence is what it does when it can.
\end{quote}

\subsection{Two constructions, and what divides them}\label{sec:compare}

It would be easy, and wrong, to rest the case on the word ``constructed''. The comparison is worth
making carefully, because three of the obvious distinctions do not hold.

\emph{Both are constructions.} Neither side hands over a closed-form solution. They exhibit data and
a solution with a property; we exhibit a flow and its outputs. That something was constructed is not
a criticism we are in a position to make.

\emph{Both depend on machines.} Theirs rests on a formalisation in Lean, ours on computation with a
posteriori error control. Neither is a proof a person has checked by hand --- and the formalised
proof runs to fifty-seven pages which, at the time of writing, no human has verified in full. We
draw no conclusion from that except the obvious symmetry: a standard admitting a machine-checked
proof that nobody has read cannot consistently reject a computation whose error is bounded by a
quantity it produces and whose predictions are checked against measurements made by other people.
The question of what it means to understand a proof is now everyone's question, and it is not an
objection one side can raise against the other.

\emph{Both choose their data.} A cylinder in a uniform stream is a choice as much as a compactly
supported vorticity field is. ``Their data is special'' is an argument only if ours is special in
some different way, and the difference has to be stated.

Three distinctions do hold.

\emph{What is certified.} They certify a property of an exact solution: with this datum, the
lifespan is finite. We certify a quantity --- that drag or lift is determined to a stated tolerance
--- together with the theorem of Section~\ref{sec:weak} that the limit solves Euler weakly. Ours is
the weaker claim about solutions and the stronger one about what can be observed.

\emph{Which function space.} Theirs is classical up to $T_*$ and conserves energy. Ours is
H\"older and dissipates. This is the substance of the section.

\emph{Where the vorticity came from.} Theirs is in the datum. Ours is made at the boundary.

\paragraph{What each method can construct at all.} There is an asymmetry running the other way,
and it should be recorded, because it bears on why a computation is used here and not as second best.

Analysis can construct dissipative H\"older solutions: that is convex integration, and
\cite{dls,bdlis,isett}. But observe what it produces --- infinitely many, for any prescribed energy
profile, including profiles along which the energy increases. The programme calls them wild and the
word is right. None is selected by anything, none issues from a configuration, and none will give
the drag of a cylinder. What is obtained is existence, not the solution.

The reason is structural. Analytical construction needs a symmetry or an ansatz to reduce the
problem, and every blowup construction buys its tractability with one: self-similarity collapses the
equations to a profile and leaves a fixed point to be established, and axisymmetry does the same
work. Turbulence affords neither. It is \emph{defined} by carrying structure at every scale with no
organising principle to exploit --- no symmetry, no similarity variable, no closed form --- and that
is not a gap in the analytical repertoire which someone may later fill, it is what the word means.
Convex integration is the exception that confirms it: it reaches dissipative solutions with no
ansatz, and pays by producing all of them at once and none in particular.

So the division is not rigorous against approximate. Analysis constructs \emph{special} solutions,
exactly, where a symmetry exists; and \emph{wild families}, where existence is wanted without
selection. Computation constructs \emph{the} solution for a given configuration, with its outputs
determined and open to check. Only the third answers the question what this flow does, and the
object it returns --- a turbulent flow past a given body, carrying a definite drag --- is not
something anyone knows how to write down in analytical steps.

\paragraph{Sought, or found.} Underneath those is a difference of direction, and we take it to be
the real one. A constructed singularity is \emph{sought}: one starts from the property wanted, looks
for data admitting it, and the work is the demonstration that the search succeeded. The conclusion
is present in the method from the first line. A computed flow is \emph{found}: one starts from a
configuration and the equations and reports what happens, and nothing in the method knows about
singularities or about turbulence.

A search ranges over the whole solution space, which is very much larger than its physical part, so
it may perfectly well return a state no fluid occupies --- and if the property sought is rare, it
will. This is the precise sense in which we call these solutions non-physical, and it is sharper
than calling them constructed. Their datum is a divergence-free field chosen so that stretching
concentrates it; it cannot be prepared, and any error in preparing it is a perturbation off the
family. A cylinder in a stream can be built, is built, in every wind tunnel, and the drag it makes
is measured by people with no stake in the answer.

We should concede that a computation can be steered too --- mesh, stabilisation and turbulence model
are all choices, and a determined computationalist finds much of what he looks for. Forwards and
backwards is a tendency of the two methods, not a guarantee about any use of them. What makes the
present case forwards is not that it is computational. It is that the configuration is one you can
build and the output is checked against a measurement nobody in the argument controls.

\paragraph{What regularity theory can reach.} It is worth adding why the question is not settled
by bounding things. The analytic tradition produces conditional results of the form: if the velocity
is bounded by $A$, then its derivatives are bounded by an expression exponential in $A$ iterated
several times. As theorems these are real. As statements about a fluid they are empty --- a bound of
size $\exp\exp\exp A$ is not one any flow can be said to respect, and no measurement could
distinguish it from no bound at all. A H\"older-$1/3$ field is neither differentiable nor unbounded,
and that is exactly the object regularity theory has no way to reach: it is not looking for a
category between smooth and singular, because the dichotomy it works in has no name for one.

\paragraph{And the setting is the point too.} The constructions are posed either on the whole space,
where by Section~\ref{sec:setting} there is no body, no slip condition and no separation --- so no
part of the mechanism of Sections~\ref{sec:dalembert}--\ref{sec:vorticity} is present --- or in a
cylinder with a boundary and an imposed symmetry, where the symmetry is what forces the flow onto
the axis and supplies the concentration. In neither case is a body being passed. The data are
special in the strong sense that they are chosen so that the self-similar profile exists, and the
proof is the verification that it does.

\paragraph{On stability.} The strongest form of the objection to the above is that blow-up has been
shown stable, and stability is what distinguishes the generic from the constructed. We would
observe that stability is always stability \emph{within} a class, and that the classes here are
symmetry classes: axisymmetric, or axisymmetric with swirl, or with a prescribed sign structure.
A perturbation that breaks the symmetry is not in the class, and it is exactly the symmetry-breaking
spanwise perturbation that in Section~\ref{sec:dalembert} is the unstable one. Stability within the
symmetric class and instability to spanwise perturbation are not in conflict; they are the same
statement seen twice, and it is the second that decides what a flow does.

\section{Summary}\label{sec:summary}

Euler flow past a body with slip has a stable attachment and an unstable separation. The separation
instability generates streamwise vorticity, the vorticity is stretched and broken down, and the
result is a cascade with an origin and no termination. The flow has bounded velocity, unbounded
gradients, H\"older exponent $1/3$, and dissipates at a rate independent of any viscosity: a
dissipative weak solution in Onsager's sense, and the dissipation is the drag. Every body separates,
so the streamlined case is the bluff case with a smaller unstable region, and there is no potential
flow regime at any Reynolds number. Without a body there is no separation and no vorticity that was
not put in by hand, which is why we take flow past a body to be the only setting in which these
equations describe anything.

\paragraph{The difference, in one line.} The singular constructions and the account above give
Euler's equations \emph{different solution concepts}, and everything else follows from that. Theirs
is conservative. A classical solution conserves energy exactly, so there is no cascade, so a growing
gradient has nowhere to go but concentration --- the singularity is not found in an unlucky corner
of the space, it is what the conservative concept leaves available. Ours is dissipative. Dissipation
requires roughness, roughness opens the cascade, the energy leaves, and nothing concentrates because
nothing has to.

At its sharpest this is the second law. A smooth Euler solution is time-reversible: run the
construction backwards and it is still a solution. A fluid is not reversible, because drag has a
sign. The dissipation anomaly is an arrow of time surviving the removal of the viscosity that was
supposed to supply it, which is what Onsager saw. So the disagreement is not about Euler's
equations, where there is nothing to disagree about. It is about whether the solutions admitted keep
an arrow of time. Smooth solutions are the reversible ones. They exist, and they are not what
happens.

\paragraph{And what each account delivers.} The trade should be stated plainly, because it runs both
ways. They prove an exact statement about an object with no observables: there exists a datum whose
smooth solution has finite lifespan --- not which data, with no measure and no genericity; not what
the flow looks like; no drag, no lift, nothing a wind tunnel could confirm or refute; and the datum
is an iterative limit, so it cannot be written down. We compute approximate statements about
observables: drag to a few per cent, the separation line, the lift curve, the stall, the drag
crisis, each with an error bound the computation itself produces and each open to check by someone
with no interest in the outcome.

Their result is mathematically substantial and physically empty; ours is physically substantial and
mathematically the weaker. Nor do we compute \emph{the} solution, and it is worth saying why not:
for weak Euler solutions uniqueness is false, by the same convex-integration programme that produced
the dissipative ones. There is no unique solution to be had. What is determined is the mean-value
outputs, and those are what G2 returns, with the limit of its sequence a weak solution of Euler by
the residual control of Section~\ref{sec:weak}.

All physical Euler flow is turbulent. The blow-up solutions are not physical --- not because they
are constructed, since so is ours, and not because the velocity becomes infinite, which it does not,
but because they conserve energy for the whole of their existence, and a fluid does not.

\begin{thebibliography}{99}

\bibitem{book} J.~Hoffman and C.~Johnson, \emph{Computational Turbulent Incompressible Flow},
Applied Mathematics: Body and Soul, vol.~4, Springer, 2007.

\bibitem{dalembert} J.~Hoffman and C.~Johnson, \emph{Resolution of d'Alembert's paradox},
J.\ Math.\ Fluid Mech.\ \textbf{12} (2010) 321--334.

\bibitem{flight} J.~Hoffman, J.~Jansson and C.~Johnson, \emph{New theory of flight},
J.\ Math.\ Fluid Mech.\ \textbf{18} (2016) 219--241.

\bibitem{drag} J.~Hoffman, \emph{Efficient computation of mean drag for the subcritical flow past a
circular cylinder using general Galerkin G2}, Int.\ J.\ Numer.\ Meth.\ Fluids \textbf{59} (2009)
1241--1258.

\bibitem{openaieuler} OpenAI, \emph{Finite time blowup for the Euler equation}, preprint, 2026,
57~pp.
\newblock \url{https://cdn.openai.com/pdf/315b36cd-ec98-4023-8342-93345194ece1/euler.pdf}

\bibitem{buckmaster} T.~Buckmaster and L.~Alp\"oge, \emph{Statement}, 2026.
\newblock \url{https://cims.nyu.edu/~tristanb/statement.pdf}

\bibitem{blowupHJ} J.~Hoffman and C.~Johnson, \emph{Blowup of incompressible Euler solutions},
BIT Numer.\ Math.\ \textbf{48} (2008) 285--307.

\bibitem{dragcrisis} J.~Hoffman, \emph{Simulation of turbulent flow past bluff bodies on coarse
meshes using General Galerkin methods: drag crisis and turbulent Euler solutions}, Comput.\ Mech.
[TO VERIFY: volume, year, pages]

\bibitem{inertialbody} \emph{Inertial momentum dissipation for viscosity solutions of Euler
equations: external flow around a smooth body}, arXiv:2206.05325. [TO VERIFY: authors]

\bibitem{kolmogorov41} A.~N.~Kolmogorov, \emph{The local structure of turbulence in incompressible
viscous fluid for very large Reynolds numbers}, Dokl.\ Akad.\ Nauk SSSR \textbf{30} (1941) 301--305.

\bibitem{taylor} G.~I.~Taylor, \emph{Statistical theory of turbulence}, Proc.\ R.\ Soc.\ Lond.\ A
\textbf{151} (1935) 421--444.

\bibitem{onsager} L.~Onsager, \emph{Statistical hydrodynamics}, Nuovo Cimento \textbf{6},
Suppl.~2 (1949) 279--287.

\bibitem{bkm} J.~T.~Beale, T.~Kato and A.~Majda, \emph{Remarks on the breakdown of smooth solutions
for the 3-D Euler equations}, Comm.\ Math.\ Phys.\ \textbf{94} (1984) 61--66.

\bibitem{cet} P.~Constantin, W.~E and E.~S.~Titi, \emph{Onsager's conjecture on the energy
conservation for solutions of Euler's equation}, Comm.\ Math.\ Phys.\ \textbf{165} (1994) 207--209.

\bibitem{eyink} G.~L.~Eyink, \emph{Energy dissipation without viscosity in ideal hydrodynamics~I.
Fourier analysis and local energy transfer}, Physica~D \textbf{78} (1994) 222--240.

\bibitem{dls} C.~De~Lellis and L.~Sz\'ekelyhidi~Jr., \emph{Dissipative Euler flows and Onsager's
conjecture}, J.\ Eur.\ Math.\ Soc.\ \textbf{16} (2014) 1467--1505.

\bibitem{bdlis} T.~Buckmaster, C.~De~Lellis, P.~Isett and L.~Sz\'ekelyhidi~Jr., \emph{Anomalous
dissipation for $1/5$-H\"older Euler flows}, Ann.\ of Math.\ \textbf{182} (2015) 127--172.

\bibitem{isett} P.~Isett, \emph{A proof of Onsager's conjecture}, Ann.\ of Math.\ \textbf{188}
(2018) 871--963.

\bibitem{luohou} G.~Luo and T.~Y.~Hou, \emph{Potentially singular solutions of the 3D axisymmetric
Euler equations}, Proc.\ Natl.\ Acad.\ Sci.\ USA \textbf{111} (2014) 12968--12973.

\bibitem{elgindi} T.~M.~Elgindi, \emph{Finite-time singularity formation for $C^{1,\alpha}$
solutions to the incompressible Euler equations on $\mathbb{R}^3$}, Ann.\ of Math.\ \textbf{194}
(2021) 647--727.

\bibitem{chenhou} J.~Chen and T.~Y.~Hou, \emph{Stable nearly self-similar blowup of the 2D
Boussinesq and 3D Euler equations with smooth data}, and \emph{Singularity formation in the 3D Euler
equations with smooth initial data and boundary}, Proc.\ Natl.\ Acad.\ Sci.\ USA (2025).
[TO VERIFY: volume and pages]

\bibitem{chen2026} J.~Chen, \emph{Asymptotically self-similar blowup for the 3D incompressible Euler
equations with $C^{1,1/3}$ velocity}, arXiv:2605.15149 (2026). [TO VERIFY]

\bibitem{stable} \emph{Stable singularity of the Euler equations on $\mathbb{R}^3$},
arXiv:2609.10867 (2026). [TO VERIFY: authors]

\end{thebibliography}

\end{document}
